% Chapter 1: Smooth Manifolds
% Lee, Introduction to Smooth Manifolds (2nd ed., GTM 218).
% Generated from the section extraction; nodes carry only Lean that is
% verified sorry-free. See ../../UPSTREAM_LEAN_AUDIT.md.


\section{Topological Manifolds}

\begin{definition}
\label{def:1-extra-2}
\lean{OpenPartialHomeomorph}

\mathlibok
A chart on a topological $n$-manifold $M$ is a pair $(U,\varphi)$, where
$U\subseteq M$ is open and $\varphi:U\to\widehat U$ is a homeomorphism onto
an open set $\widehat U\subseteq\mathbb R^n$. It is centered at $p\in U$ when
$\varphi(p)=0$; replacing $\varphi$ by $\varphi-\varphi(p)$ centers any chart
at $p$. The set $U$ is a coordinate domain, and it is a coordinate ball or
coordinate cube when $\widehat U$ is respectively an open ball or open cube.
Writing
\[
\varphi(p)=(x^1(p),\ldots,x^n(p))
\]
defines the local coordinates $x^1,\ldots,x^n$ on $U$.
\end{definition}

\begin{example}[Graphs of Continuous Functions]
\label{ex:1-3}
\lean{graph_coordinates}
\leanok


\dcref{ch1:1.3}
For an open set $U\subseteq\mathbb R^n$ and a continuous map
$f:U\to\mathbb R^k$, let
\[
\Gamma(f)=\{(x,y)\in\mathbb R^n\times\mathbb R^k:x\in U,\ y=f(x)\}.
\]
The restriction $\varphi=\pi_1|_{\Gamma(f)}:\Gamma(f)\to U$ is a
homeomorphism with inverse $x\mapsto(x,f(x))$. Thus $\Gamma(f)$ is a
topological $n$-manifold with the global chart $(\Gamma(f),\varphi)$, called
graph coordinates. The same construction applies when any $k$ coordinates
are continuous functions of the remaining $n$ coordinates in an open subset
of $\mathbb R^n$.
\end{example}

\begin{example}[Projective Spaces]
\label{ex:1-5}
\lean{real_projective_space_has_standard_chart}
\leanok


\dcref{ch1:1.5}
Define $\mathbb{RP}^n$ as the set of one-dimensional subspaces of
$\mathbb R^{n+1}$ with the quotient topology induced by
\[
\pi:\mathbb R^{n+1}\smallsetminus\{0\}\to\mathbb{RP}^n,
\qquad x\mapsto[x].
\]
For $1\leq i\leq n+1$, set
$\widetilde U_i=\{x:x^i\neq0\}$ and $U_i=\pi(\widetilde U_i)$. These are open,
because $\widetilde U_i$ is saturated and open, and they cover
$\mathbb{RP}^n$. The quotient property defines continuous maps
\[
\varphi_i[x^1,\ldots,x^{n+1}]
=\left(\frac{x^1}{x^i},\ldots,\frac{x^{i-1}}{x^i},
\frac{x^{i+1}}{x^i},\ldots,\frac{x^{n+1}}{x^i}\right).
\]
Each $\varphi_i:U_i\to\mathbb R^n$ is a homeomorphism, with inverse
\[
\varphi_i^{-1}(u^1,\ldots,u^n)
=[u^1,\ldots,u^{i-1},1,u^i,\ldots,u^n].
\]
Hence $\mathbb{RP}^n$ is locally Euclidean of dimension $n$.
\end{example}

\begin{proposition}[Topological Manifold Structure of Projective Space]
\label{exer:1-6}
\lean{realProjectiveSpace_chartDomain_isOpen}
\leanok


\dcref{ch1:1.6}
$\mathbb{RP}^n$ is Hausdorff and second-countable, and hence is a topological
$n$-manifold.
\end{proposition}

\begin{proposition}[Compactness of Projective Space]
\label{exer:1-7}
\lean{realProjectiveSpaceCompactSpace}


\dcref{ch1:1.7}
$\mathbb{RP}^n$ is compact; the restriction
$\pi|_{\mathbb S^n}:\mathbb S^n\to\mathbb{RP}^n$ is surjective.
\end{proposition}

\begin{example}[Product Manifolds]
\label{ex:1-8}
\lean{isManifold_pi}
\leanok


\dcref{ch1:1.8}
If $M_i$ is a topological manifold of dimension $n_i$ for $1\leq i\leq k$,
then $M_1\times\cdots\times M_k$ is a topological manifold of dimension
$n_1+\cdots+n_k$. Indeed, products preserve the Hausdorff and
second-countability properties, and charts $(U_i,\varphi_i)$ yield the product
chart
\[
\varphi_1\times\cdots\times\varphi_k:
U_1\times\cdots\times U_k\longrightarrow\mathbb R^{n_1+\cdots+n_k}.
\]
\end{example}

\begin{example}[Tori]
\label{ex:1-9}
\lean{n_torus}
\leanok


\dcref{ch1:1.9}
For $n\geq1$, the $n$-torus
$\mathbb T^n=(\mathbb S^1)^n$ is a topological $n$-manifold.
\end{example}

\begin{definition}
\label{def:1-extra-3}
\lean{LocPathConnectedSpace}

\mathlibok
A topological space $X$ is connected if it is not the union of two disjoint
nonempty open sets, path-connected if every two points are joined by a path,
and locally path-connected if it has a basis of path-connected open sets.
\end{definition}

\begin{definition}
\label{def:1-extra-4}
\lean{TopologicalSpace.IsOpenCover}

\mathlibok
A family $\mathcal X$ of subsets of a topological space $M$ is locally finite
if every point has a neighborhood meeting only finitely many members of
$\mathcal X$. A cover $\mathcal V$ refines a cover $\mathcal U$ if every
$V\in\mathcal V$ is contained in some $U\in\mathcal U$. The space $M$ is
paracompact if every open cover has an open locally finite refinement.
\end{definition}

\begin{lemma}
\label{lem:1-13}
\lean{LocallyFinite.closure_iUnion}

\mathlibok
\dcref{ch1:1.13}
Let $\mathcal X$ be a locally finite family of subsets of a topological space
$M$. Then $\{\overline X:X\in\mathcal X\}$ is locally finite and
\[
\overline{\bigcup_{X\in\mathcal X}X}
=\bigcup_{X\in\mathcal X}\overline X.
\]
\end{lemma}



\section{Smooth Structures}

\begin{definition}
\label{def:1-2-extra-1}
\lean{Diffeomorph.toHomeomorph}

\mathlibok
For open sets $U\subseteq\mathbb R^n$ and $V\subseteq\mathbb R^m$, a map
$F:U\to V$ is smooth if all component functions have continuous partial
derivatives of every order. A bijective smooth map with smooth inverse is a
diffeomorphism, and every diffeomorphism is a homeomorphism.
\end{definition}

\begin{definition}
\label{def:1-2-extra-2}
\lean{e.symm}

\mathlibok
For charts $(U,\varphi)$ and $(V,\psi)$ on a topological $n$-manifold, the
transition map on a nonempty overlap is
\[
\psi\circ\varphi^{-1}:\varphi(U\cap V)\longrightarrow\psi(U\cap V).
\]
The charts are smoothly compatible if $U\cap V=\varnothing$ or this map is a
diffeomorphism between the indicated open subsets of $\mathbb R^n$.
\end{definition}

\begin{definition}
\label{def:1-2-extra-3}
\lean{OpenPartialHomeomorph.IsSmoothAtlas}
\leanok


An atlas on $M$ is a family of charts whose domains cover $M$. It is smooth
when every pair of its charts is smoothly compatible.
\end{definition}

\begin{definition}
\label{def:1-2-extra-4}
\lean{StructureGroupoid.maximalAtlas}

\mathlibok
A smooth atlas $\mathcal A$ is maximal if it is not properly contained in
another smooth atlas; equivalently, every chart smoothly compatible with all
charts in $\mathcal A$ already belongs to $\mathcal A$.
\end{definition}

\begin{definition}
\label{def:1-2-extra-5}
\lean{smooth_manifold_smooth_structure}
\leanok


A smooth structure on a topological manifold is a maximal smooth atlas. A
smooth manifold is a topological manifold equipped with a smooth structure.
A smooth manifold structure consists of both the manifold topology and the
smooth structure.
\end{definition}

\begin{proposition}
\label{prop:1-17}
\lean{same_smooth_structure_iff_union_is_smooth_atlas}
\leanok


\dcref{ch1:1.17}
Let $M$ be a topological manifold.
\begin{enumerate}
\item Every smooth atlas $\mathcal A$ on $M$ is contained in a unique maximal
smooth atlas, called the smooth structure determined by $\mathcal A$.
\item Two smooth atlases determine the same smooth structure if and only if
their union is a smooth atlas.
\end{enumerate}
\end{proposition}

\begin{proof}
\leanok
\sketch Let $\overline{\mathcal A}$ be the family of charts compatible with
every member of $\mathcal A$. Given $(U,\varphi),(V,\psi)\in
\overline{\mathcal A}$ and $p\in U\cap V$, choose $(W,\theta)\in\mathcal A$
with $p\in W$. Near $\varphi(p)$,
\[
\psi\circ\varphi^{-1}
=(\psi\circ\theta^{-1})\circ(\theta\circ\varphi^{-1}),
\]
so the two charts are compatible. Thus $\overline{\mathcal A}$ is smooth and
is maximal by its definition. Every maximal atlas containing $\mathcal A$ lies
in $\overline{\mathcal A}$ and therefore equals it. For the second assertion,
a smooth union lies in its unique maximal extension; conversely, two atlases
inside one maximal atlas have a smooth union.
\end{proof}


\begin{remark}
\label{rem:1-2-extra-6}
\lean{IsManifold}

\mathlibok
Requiring every transition map and its inverse to be $C^k$ gives a $C^k$
structure; $C^0$ structures are topological manifold structures. Requiring
both directions to be real-analytic gives a real-analytic, or $C^\omega$,
structure. On an even-dimensional manifold, identifying $\mathbb R^{2m}$ with
$\mathbb C^m$ and requiring both directions to be complex-analytic gives a
complex manifold structure.
\end{remark}


\section{Local Coordinate Representations}

\begin{definition}
\label{def:1-3-extra-1}
\lean{IsRegularCoordinateBall}
\leanok


A chart in the maximal smooth atlas is a smooth chart; its domain is a smooth
coordinate domain. Smooth coordinate balls and cubes have Euclidean balls and
cubes as their respective coordinate images. A subset $B\subseteq M$ is a
regular coordinate ball if there are a smooth coordinate ball $B'$ containing
$\overline B$, a smooth coordinate map $\varphi:B'\to\mathbb R^n$, and
$0<r<r'$ such that
\[
\varphi(B)=B_r(0),\qquad
\varphi(\overline B)=\overline B_r(0),\qquad
\varphi(B')=B_{r'}(0).
\]
\end{definition}

\begin{proposition}
\label{prop:1-19}
\lean{exists_countable_regular_coordinate_ball_basis}
\leanok


\dcref{ch1:1.19}
Every smooth manifold has a countable basis of regular coordinate balls.
\end{proposition}


\begin{notation}
\label{not:1-3-extra-2}
\lean{extChartAt}

\mathlibok
In a chart $(U,\varphi)$ with coordinate functions $(x^1,\ldots,x^n)$, the
local coordinate representation of $p\in U$ is
$\varphi(p)=(x^1(p),\ldots,x^n(p))$; one may write
$p=(x^1,\ldots,x^n)$ in these coordinates.
\end{notation}


\begin{example}[Euclidean Spaces]
\label{ex:1-22}
\lean{euclideanSpace_standard_smooth_structure}
\leanok


\dcref{ch1:1.22}
The standard smooth structure on $\mathbb R^n$ is determined by
$(\mathbb R^n,\operatorname{Id}_{\mathbb R^n})$. Its smooth charts are exactly
the pairs $(U,\varphi)$ for which $\varphi$ is a Euclidean diffeomorphism from
$U$ onto an open subset of $\mathbb R^n$. This is the structure used unless
another is specified.
\end{example}

\begin{example}[Another Smooth Structure on $\mathbb{R}$]
\label{ex:1-23}
\lean{cubicChart_not_mem_standard_smooth_maximalAtlas}


\dcref{ch1:1.23}
The homeomorphism
\[
\psi:\mathbb R\to\mathbb R,\qquad \psi(x)=x^3 \tag{1.1}
\]
determines a smooth structure using the one-chart atlas
$(\mathbb R,\psi)$. It differs from the standard structure because the
transition $\psi^{-1}(y)=y^{1/3}$ is not smooth at $0$.
More generally, every positive-dimensional topological manifold that admits
one smooth structure admits uncountably many distinct smooth structures.
\end{example}

\begin{example}[Finite-Dimensional Vector Spaces]
\label{ex:1-24}
\lean{basis_coordinate_change_contDiff}
\leanok


\dcref{ch1:1.24}
Let $V$ be an $n$-dimensional real vector space with its norm-independent
topology. An ordered basis $(E_1,\ldots,E_n)$ defines
\[
E:\mathbb R^n\to V,\qquad E(x)=\sum_{i=1}^n x^iE_i,
\]
and hence a chart $(V,E^{-1})$. For another basis $\widetilde E_j$ with
$E_i=\sum_j A_i^j\widetilde E_j$, the transition map is the invertible linear
map
\[
\widetilde x^j=\sum_i A_i^j x^i.
\]
Therefore all basis charts are smoothly compatible and determine the standard
smooth structure on $V$.
\end{example}

\begin{notation}
\label{not:1-4-extra-1}
\lean{Basis.sum_repr}

\mathlibok
The Einstein convention interprets an index occurring exactly twice in a
monomial, once upstairs and once downstairs, as summed over its full range;
thus $x^iE_i=\sum_{i=1}^n x^iE_i$. Repeated indices are dummy indices. Basis
vectors carry lower indices, while vector components and Euclidean coordinates
carry upper indices. Summations violating this upper--lower pattern, including
the Euclidean dot product $x\cdot y=\sum_i x^iy^i$ and formulas involving the
standard symplectic form on $\mathbb R^{2n}$, are written explicitly.
\end{notation}

\begin{example}[The General Linear Group]
\label{ex:1-27}
\lean{Units.isOpenEmbedding_val}

\mathlibok
\dcref{ch1:1.27}
The group $\operatorname{GL}(n,\mathbb R)$ is the nonvanishing locus of the
continuous determinant on the $n^2$-dimensional vector space of real
$n\times n$ matrices. It is therefore an open smooth submanifold of dimension
$n^2$.
\end{example}

\begin{example}[Matrices of Full Rank]
\label{ex:1-28}
\lean{fullRankRange_iff_natLeRank}
\leanok


\dcref{ch1:1.28}
If $m<n$, the rank-$m$ matrices in
$\operatorname M(m\times n,\mathbb R)$ form an open set: at each such matrix,
some $m\times m$ minor is nonzero, and remains so nearby. Hence this full-rank
locus is a smooth $mn$-manifold. The analogous statement holds for rank-$n$
matrices when $n<m$.
\end{example}

\begin{example}[Graphs of Smooth Functions]
\label{ex:1-30}
\lean{graph_smooth_structure}
\uses{ex:1-3}
\leanok



\dcref{ch1:1.30}
If $U\subseteq\mathbb R^n$ is open and $f:U\to\mathbb R^k$ is smooth, the
graph $\Gamma(f)$ has the canonical smooth structure determined by the global
graph-coordinate chart $\pi_1|_{\Gamma(f)}:\Gamma(f)\to U$.
\end{example}

\begin{example}[Projective Spaces]
\label{ex:1-33}
\lean{realProjectiveChartTransitionDiffeomorph}
\uses{ex:1-5}
\leanok



\dcref{ch1:1.33}
The charts $(U_i,\varphi_i)$ on $\mathbb{RP}^n$ are smoothly compatible. For
$i>j$, their transition is
\[
\varphi_j\circ\varphi_i^{-1}(u^1,\ldots,u^n)
=\left(\frac{u^1}{u^j},\ldots,\frac{u^{j-1}}{u^j},
\frac{u^{j+1}}{u^j},\ldots,\frac{u^{i-1}}{u^j},
\frac1{u^j},\frac{u^i}{u^j},\ldots,\frac{u^n}{u^j}\right),
\]
a diffeomorphism between the corresponding overlap images.
\end{example}

\begin{proposition}[Angle Functions and Circle Charts]
\label{prob:1-8}
\lean{IsAngleFunction}
\leanok


For an open set $U\subseteq\mathbb S^1\subseteq\mathbb C$, an
\emph{angle function} is a continuous map $\theta:U\to\mathbb R$ satisfying
$e^{i\theta(z)}=z$. Such a function exists if and only if
$U\neq\mathbb S^1$, and every angle function is a smooth chart for the
standard smooth structure on $\mathbb S^1$.
\end{proposition}

\begin{lemma}[Smooth Manifold Chart Lemma]
\label{lem:1-35}
\lean{ChartedSpaceCore.eq_toTopologicalSpace_of_chartedSpace}
\leanok


\dcref{ch1:1.35}
Let $M$ be a set with subsets $U_\alpha$ and maps
$\varphi_\alpha:U_\alpha\to\mathbb R^n$ satisfying:
\begin{enumerate}
\item each $\varphi_\alpha$ bijects $U_\alpha$ with an open subset of
$\mathbb R^n$;
\item both $\varphi_\alpha(U_\alpha\cap U_\beta)$ and
$\varphi_\beta(U_\alpha\cap U_\beta)$ are open;
\item on every nonempty overlap,
$\varphi_\beta\circ\varphi_\alpha^{-1}$ is smooth;
\item a countable subfamily of the $U_\alpha$ covers $M$;
\item for distinct $p,q\in M$, either some $U_\alpha$ contains both, or there
are disjoint $U_\alpha,U_\beta$ containing $p,q$, respectively.
\end{enumerate}
Then $M$ has a unique smooth manifold structure for which all
$(U_\alpha,\varphi_\alpha)$ are smooth charts.
\end{lemma}

\begin{proof}
\leanok
\sketch Use the sets $\varphi_\alpha^{-1}(V)$, with
$V\subseteq\mathbb R^n$ open, as a basis. Conditions (ii)--(iii) give
\[
\varphi_\alpha^{-1}(V)\cap\varphi_\beta^{-1}(W)
=\varphi_\alpha^{-1}\!\left(
V\cap(\varphi_\beta\circ\varphi_\alpha^{-1})^{-1}(W)\right),
\]
so basis intersections are open. Condition (i) makes the resulting space
locally Euclidean, (iv) gives second-countability, and (v) gives the Hausdorff
property. Condition (iii) then supplies a smooth atlas. The required topology
and maximal smooth atlas are forced by the given charts, proving uniqueness.
\end{proof}

\begin{example}[Grassmann Manifolds]
\label{ex:1-36}
\lean{grassmannian.chart_equiv}
\uses{ex:1-28, lem:1-35}
\leanok



\dcref{ch1:1.36}
Let $V$ be an $n$-dimensional real vector space and $0\leq k\leq n$. The set
$\operatorname G_k(V)$ of $k$-dimensional subspaces has a natural smooth
structure of dimension $k(n-k)$; in particular,
$\operatorname G_1(\mathbb R^{n+1})=\mathbb{RP}^n$.

For a decomposition $V=P\oplus Q$ with $\dim P=k$, define
\[
U_Q=\{S\in\operatorname G_k(V):S\cap Q=\{0\}\}.
\]
Every $S\in U_Q$ is the graph of the unique map
\[
X=(\pi_Q|_S)\circ(\pi_P|_S)^{-1}:P\to Q,
\]
so $S\mapsto X$ bijects $U_Q$ with
$\operatorname L(P;Q)\cong\mathbb R^{k(n-k)}$.

For a second decomposition $V=P'\oplus Q'$, write
\[
A=\pi_{P'}|_P,\qquad B=\pi_{Q'}|_P,\qquad
C=\pi_{P'}|_Q,\qquad D=\pi_{Q'}|_Q.
\]
On the overlap, $A+CX$ has full rank, so the overlap is open by
Example~\ref{ex:1-28}; the coordinate change is
\[
X'=(B+DX)(A+CX)^{-1}.
\]
Its matrix entries are smooth because matrix inversion has rational entries on
the invertible locus. A fixed basis of $V$ yields finitely many such domains
$U_Q$ covering $\operatorname G_k(V)$. Finally, any two $k$-planes have a
common complementary $(n-k)$-plane, so they lie in one chart domain. The chart
lemma therefore gives the asserted smooth manifold structure.
\end{example}


\section{Manifolds with Boundary}

\begin{definition}
\label{def:1-5-extra-1}
\lean{TopologicalManifoldWithBoundary}
\leanok


The closed upper half-space is
\[
\mathbb H^n=\{(x^1,\ldots,x^n)\in\mathbb R^n:x^n\geq0\}.
\]
For $n>0$,
\[
\operatorname{Int}\mathbb H^n=\{x:x^n>0\},\qquad
\partial\mathbb H^n=\{x:x^n=0\};
\]
for $n=0$, $\mathbb H^0=\operatorname{Int}\mathbb H^0=\mathbb R^0$ and
$\partial\mathbb H^0=\varnothing$.

A topological $n$-manifold with boundary is a Hausdorff second-countable space
locally homeomorphic to an open subset of $\mathbb R^n$ or to a relatively
open subset of $\mathbb H^n$. A chart is a homeomorphism
$\varphi:U\to W$ from an open set $U$ in the manifold onto such an open or
relatively open set $W$. It is an interior chart when $W$ is open in
$\mathbb R^n$, and a boundary chart when $W$ is relatively open in
$\mathbb H^n$ and meets $\partial\mathbb H^n$. A boundary chart with image
$B_r(x)\cap\mathbb H^n$, where $x\in\partial\mathbb H^n$, is a coordinate
half-ball.

A point is interior if it lies in an interior-chart domain. It is a boundary
point if some boundary chart sends it to $\partial\mathbb H^n$. These sets are
denoted $\operatorname{Int}M$ and $\partial M$. Every point belongs to at least
one of them: a point represented away from $\partial\mathbb H^n$ has an
interior chart after restriction.
\end{definition}

\begin{theorem}[Topological Invariance of the Boundary]
\label{thm:1-37}
\lean{ModelWithCorners.disjoint_interior_boundary}

\mathlibok
\dcref{ch1:1.37}
If $M$ is a topological manifold with boundary, then every point lies in
exactly one of $\partial M$ and $\operatorname{Int}M$. Thus these sets are
disjoint and cover $M$.
\end{theorem}


\section{Smooth Structures on Manifolds with Boundary}

\begin{definition}
\label{def:1-6-extra-2}
\lean{SmoothManifoldWithBoundary}
\leanok


A smooth structure on a topological manifold with boundary is a maximal atlas
whose transition maps and inverses are smooth. A manifold with boundary
equipped with such a structure is a smooth manifold with boundary. Every
smooth manifold is a smooth manifold with empty boundary.
\end{definition}

\begin{definition}
\label{def:1-6-extra-3}
\lean{IsRegularCoordinateHalfBall}
\leanok


A chart in the maximal atlas of a smooth manifold with boundary is smooth.
A smooth coordinate ball or half-ball is a smooth chart domain whose image is,
respectively, a Euclidean ball or a set $B_r(x)\cap\mathbb H^n$ with
$x\in\partial\mathbb H^n$. Regular coordinate balls are defined in
Definition~\ref{def:1-3-extra-1}. A subset $B\subseteq M$ is a regular
coordinate half-ball if there are a smooth coordinate half-ball $B'$
containing $\overline B$, a smooth map $\varphi:B'\to\mathbb H^n$, and
$0<r<r'$ such that
\[
\varphi(B)=B_r(0)\cap\mathbb H^n,\qquad
\varphi(\overline B)=\overline B_r(0)\cap\mathbb H^n,\qquad
\varphi(B')=B_{r'}(0)\cap\mathbb H^n.
\]
\end{definition}

\begin{proposition}[Open Subsets of Manifolds with Boundary]
\label{exer:1-44}
\lean{manifoldInterior_boundaryless}
\leanok


\dcref{ch1:1.44}
Let $M$ be a smooth $n$-manifold with boundary and let $U\subseteq M$ be open.
Then:
\begin{enumerate}
\item $U$ is a topological $n$-manifold with boundary, with smooth structure
given by the smooth charts of $M$ whose domains lie in $U$;
\item if $U\subseteq\operatorname{Int}M$, then $U$ is a smooth manifold
without boundary;
\item $\operatorname{Int}M$ is an open smooth submanifold of $M$ without
boundary.
\end{enumerate}
\end{proposition}

\begin{proposition}
\label{prop:1-45}
\lean{ModelWithCorners.boundary_of_boundaryless_left}

\mathlibok
\dcref{ch1:1.45}
If $M_1,\ldots,M_k$ are smooth manifolds and $N$ is a smooth manifold with
boundary, then $M_1\times\cdots\times M_k\times N$ is a smooth manifold with
boundary and
\[
\partial(M_1\times\cdots\times M_k\times N)
=M_1\times\cdots\times M_k\times\partial N.
\]
\end{proposition}

\begin{theorem}[Smooth Invariance of the Boundary]
\label{thm:1-46}
\lean{smooth_boundary_chart_frontier_independence}
\leanok


\dcref{ch1:1.46}
Let $M$ be a smooth manifold with boundary and $p\in M$. If one smooth chart
maps $p$ to $\partial\mathbb H^n$, then every smooth chart containing $p$ does
so.
\end{theorem}

\begin{proof}
\leanok
\sketch Suppose an interior chart $(U,\psi)$ and a boundary chart
$(V,\varphi)$ both contain $p$, with $\varphi(p)\in\partial\mathbb H^n$. Put
$\tau=\varphi\circ\psi^{-1}$, $x_0=\psi(p)$, and $y_0=\varphi(p)$. Smooth
compatibility extends $\tau^{-1}$ near $y_0$ to a smooth map $\eta$. On a
small Euclidean ball $B$ about $x_0$, one has
$\eta\circ\tau=\operatorname{Id}$, hence
\[
D\eta(\tau(x))\circ D\tau(x)=\operatorname{Id}\qquad(x\in B).
\]
Thus $D\tau$ is nonsingular on $B$, and the inverse function theorem implies
that $\tau(B)$ is open in $\mathbb R^n$. But $\tau(B)$ contains $y_0$ and lies
in $\mathbb H^n$, contradicting $y_0\in\partial\mathbb H^n$.
\end{proof}

\begin{exercise}
\label{exer:1-14}
\lean{LocallyFinite.closure_iUnion}
\mathlibok
\uses{lem:1-13}
\dcref{ch1:1.14}
\begin{itemize}
\item Exercise 1.14. Prove the preceding lemma.
\end{itemize}
\end{exercise}

\begin{exercise}
\label{exer:1-18}
\lean{same_smooth_structure_iff_union_is_smooth_atlas}
\leanok
\uses{prop:1-17}
\dcref{ch1:1.18}
\begin{itemize}
\item Exercise 1.18. Prove Proposition 1.17(b).
\end{itemize}
\end{exercise}

\begin{exercise}
\label{exer:1-20}
\lean{exists_countable_regular_coordinate_ball_basis}
\leanok
\dcref{ch1:1.20}
Prove Proposition 1.19.
\end{exercise}
